\subsubsection{真实输入下的三维隐含属性统计}\label{app:real-input-40-arity-3d}
在 40 状态核上引入三维观测
\[
\varphi_e=e,\qquad
\varphi_-=\mathbf{1}_{\{t\in Q_-\}},\qquad
\varphi_2=\mathbf{1}_{\{d=2\}},
\]
并定义加权矩阵
\[
(M_\theta)_{ij}=\sum_{i\to j}\exp(\theta_e\varphi_e+\theta_-\varphi_-+\theta_2\varphi_2),
\qquad
P(\theta)=\log\rho(M_\theta).
\]
在 $\theta=0$ 处，梯度（典型密度）为
\[
\partial_{\theta_e}P(0)=\frac{23-7\sqrt5}{20}\ (\approx 0.3673762079),\quad
\partial_{\theta_-}P(0)=\frac{5-\sqrt5}{20}\ (\approx 0.1381966011),\quad
\partial_{\theta_2}P(0)=\frac{3-\sqrt5}{10}\ (\approx 0.0763932023),
\]
Hessian（协方差密度）数值为
\[
\nabla^2P(0)\approx
\begin{pmatrix}
0.048741 & 0.032574 & 0.015964\\
0.032574 & 0.275035 & 0.115735\\
0.015964 & 0.115735 & 0.067331
\end{pmatrix},
\qquad
\bigl(\nabla^2P(0)\bigr)^{-1}\approx
\begin{pmatrix}
22.4165 & -1.5120 & -2.7160\\
-1.5120 & 13.2430 & -22.4049\\
-2.7160 & -22.4049 & 54.0075
\end{pmatrix}.
\]
因此局部大偏差率函数满足
\[
I(\alpha)\approx \tfrac12(\alpha-\alpha_\star)^\top\bigl(\nabla^2P(0)\bigr)^{-1}(\alpha-\alpha_\star),
\qquad
\alpha_\star=\left(\frac{23-7\sqrt5}{20},\frac{5-\sqrt5}{20},\frac{3-\sqrt5}{10}\right).
\]
数值复现见脚本 \texttt{scripts/exp\_sync\_kernel\_real\_input\_40.py}。

\begin{proposition}[碰撞势的三次压力方程]\label{prop:real-input-40-collision-pressure}
令 $X$ 为黄金均值移位，取其邻接矩阵
\(
F=\begin{pmatrix}1&1\\[2pt]1&0\end{pmatrix}
\)，则 $X\times X$ 的邻接矩阵为 $A:=F\otimes F$（按状态 $00,01,10,11$ 排列）。
对碰撞势 $\varphi_2=\mathbf{1}_{\{11\}}$，令 $u=e^\theta$ 并定义权矩阵
\[
W(u):=D(u)A,\qquad D(u):=\mathrm{diag}(1,1,1,u).
\]
则特征多项式分解为
\[
\det(\lambda I-W(u))=(\lambda+1)\Bigl(\lambda^3-2\lambda^2-(u+1)\lambda+u\Bigr),
\]
从而 Perron 根 $\rho(u)$ 为三次方程
\[
\rho^3-2\rho^2-(u+1)\rho+u=0
\]
的最大实根，并有
\[
P_2(\theta)=\log\rho(e^\theta).
\]
\end{proposition}

\begin{corollary}[碰撞势均值与方差的 $\sqrt5$ 闭式]\label{cor:real-input-40-collision-cumulants}
在 $u=1$（即 $\theta=0$，$\rho(1)=\lambda=\varphi^2$）处，
\[
P_2'(0)=\frac{3-\sqrt5}{10},
\qquad
P_2''(0)=\frac{6\sqrt5-5}{125}.
\]
其中 $P_2''(0)$ 与上式 Hessian 的第三对角元一致。
\end{corollary}

\begin{corollary}[三、四阶累积量闭式]\label{cor:real-input-40-collision-cumulants-higher}
在同样设定下，碰撞势压力的三、四阶导数在 $\theta=0$ 处亦有 $\sqrt5$ 闭式：
\[
\boxed{\;
P_2^{(3)}(0)=\frac{9+10\sqrt5}{625},
\qquad
P_2^{(4)}(0)=\frac{7+24\sqrt5}{3125}.
\;}
\]
\end{corollary}

\begin{corollary}[单位根扭曲下的四阶谱隙渐近]\label{cor:real-input-40-collision-twist-fourth}
令 $t\in\RR$，取单位根扭曲 $u=e^{it}$ 并记
\(
\rho(t):=\rho(e^{it})
\)。
则当 $t\to 0$ 时有
\[
\Re\bigl(P_2(it)-P_2(0)\bigr)
=-\frac{P_2''(0)}{2}\,t^2+\frac{P_2^{(4)}(0)}{24}\,t^4+O(t^6),
\]
从而（记 $\lambda=\rho(1)=\varphi^2$）
\[
\boxed{\
\log\frac{\rho(t)}{\lambda}
=-\frac{6\sqrt5-5}{250}\,t^2+\frac{7+24\sqrt5}{75000}\,t^4+O(t^6).
\ }
\]
特别地，对奇素数模 $p$ 的最小非平凡角 $t_p:=2\pi/p$，
\[
\log\frac{\rho(t_p)}{\lambda}
=-\frac{6\sqrt5-5}{250}\left(\frac{2\pi}{p}\right)^2
\frac{7+24\sqrt5}{75000}\left(\frac{2\pi}{p}\right)^4
O\!\left(p^{-6}\right).
\]
\end{corollary}

\begin{remark}[模 $2$ 的碰撞奇偶扭曲：极简三次代数数]\label{rem:real-input-40-collision-parity}
当只记录“碰撞次数的奇偶”时取 $u=-1$，三次方程 \eqref{prop:real-input-40-collision-pressure} 中 $(u+1)$ 项消失，得到
\[
\rho^3-2\rho^2-1=0.
\]
因此非平凡谱半径 $\rho(-1)$ 即为该三次方程的最大实根（数值上 $\rho(-1)\approx 2.2055694304$，从而 $\rho(-1)/\lambda\approx 0.8425$），可作为“near-coboundary 慢混合”现象的一个对照基准：模 $2$ 扭曲远离 coboundary，谱隙更大、混合更快。
\end{remark}

\begin{remark}[碰撞频率的大偏差参数化]\label{rem:real-input-40-collision-ldp-param}
由三次方程可消去 $u$，得到参数化表示
\[
u=\frac{\rho(\rho^2-2\rho-1)}{\rho-1}\qquad(\rho>1).
\]
进一步记 $\alpha(\theta)=P_2'(\theta)$，则可写成
\[
\alpha(\rho)=\frac{(\rho^2-2\rho-1)(\rho-1)}{2\rho^3-5\rho^2+4\rho+1},
\]
从而碰撞频率的速率函数可参数化为
\[
I_2(\alpha(\rho))=\alpha(\rho)\log u(\rho)-\log\rho,\qquad \rho\in(1,\infty).
\]
\end{remark}

\begin{proposition}[全实谱与根型在 \texorpdfstring{$u>0$}{u>0} 上固定]\label{prop:real-input-40-collision-all-real}
对任意 \(u>0\)，三次多项式
\[
f(\lambda;u):=\lambda^3-2\lambda^2-(u+1)\lambda+u
\]
的判别式为
\[
\Delta(u)=4u^3+25u^2+88u+8>0.
\]
因此 \(f(\lambda;u)\) 在 \(\RR\) 上恒有三个互异实根；结合额外特征值 \(-1\)，矩阵 \(W(u)=D(u)A\) 的全部特征值均为实数。
并且根型（按大小/符号）满足：
\[
\lambda_+(u)>1,\qquad 0<\lambda_0(u)<1,\qquad \lambda_-(u)<0,
\]
其中 \(\lambda_+(u)\) 为 Perron 根（命题 \ref{prop:real-input-40-collision-pressure}）。
更具体地，\(u\le 1\) 时有 \(\lambda_-(u)\in[-1,0)\)，而 \(u>1\) 时有 \(\lambda_-(u)<-1\)。
\end{proposition}
\begin{proof}
判别式直接计算得 \(\Delta(u)=4u^3+25u^2+88u+8>0\)（\(u>0\) 显然），故三根互异且全实。
又有 \(f(0;u)=u>0,\ f(1;u)=1-2-(u+1)+u=-2<0\)，故在 \((0,1)\) 内存在一根 \(\lambda_0(u)\)。
并且 \(f(\lambda;u)\to+\infty\) 当 \(\lambda\to+\infty\)，结合 \(f(1;u)<0\) 得在 \((1,\infty)\) 内存在一根 \(\lambda_+(u)>1\)。
由 Vieta 关系，三根乘积为 \(\lambda_+(u)\lambda_0(u)\lambda_-(u)=-u<0\)，故第三根 \(\lambda_-(u)\) 为负数。
最后计算
\(
f(-1;u)=2(u-1)
\)：
当 \(u\le 1\) 时 \(f(-1;u)\le 0\) 且 \(f(0;u)>0\)，故负根落在 \([-1,0)\)；
当 \(u>1\) 时 \(f(-1;u)>0\) 且 \(f(\lambda;u)\to-\infty\) 当 \(\lambda\to-\infty\)，故负根满足 \(\lambda_-(u)<-1\)。
证毕。
\end{proof}

\begin{proposition}[碰撞位势的核版 RH/Ramanujan 窗口与临界五次方程]\label{prop:real-input-40-collision-rhk-window}
令 \(\lambda_1(u):=\lambda_+(u)\) 为 Perron 根，并记次谱半径
\[
\Lambda(u):=\max\{|\lambda_0(u)|,|\lambda_-(u)|,1\}.
\]
则核版 RH 条件 \(\textsf{RH}(u)\)（定义 \ref{def:kernel-rh}）在该 \(4\times4\) 碰撞位势上等价于
\[
\Lambda(u)\le \lambda_1(u)^{1/2}.
\]
存在唯一的临界参数 \(u_R>1\) 使得
\[
u_R^5+4u_R^4+3u_R^3-96u_R^2-42u_R-14=0,
\]
并且
\[
\textsf{RH}(u)\ \Longleftrightarrow\ 0<u\le u_R.
\]
数值上 \(u_R\approx 3.59262181344\)（\(\theta_R=\log u_R\approx 1.27888\)）。
\end{proposition}
\begin{proof}
由命题 \ref{prop:real-input-40-collision-all-real}，当 \(u\le 1\) 时 \(|\lambda_-(u)|\le 1\) 且 \(0<\lambda_0(u)<1\)，故 \(\Lambda(u)=1\)，而 \(\lambda_1(u)>1\) 给出 \(1\le \sqrt{\lambda_1(u)}\)，于是 \(\textsf{RH}(u)\) 自动成立。

当 \(u>1\) 时有 \(\lambda_-(u)<-1\)，故 \(\Lambda(u)=-\lambda_-(u)\)。
因此 \(\textsf{RH}(u)\) 等价于
\[
(-\lambda_-(u))^2\le \lambda_+(u).
\]
令 \(b:=-\lambda_-(u)>1\)。临界等号发生当且仅当存在 \(b>1\) 使得
\(
f(-b;u)=0
\)
且
\(
f(b^2;u)=0
\)。
消去参数 \(b\)（对 \(b\) 取 resultant）得到关于 \(u\) 的代数方程
\(
u(u^5+4u^4+3u^3-96u^2-42u-14)=0
\)；
由于 \(u>0\)，临界点由五次多项式的正实根给出。该五次的系数符号序列为 \((+,+,+,-,-,-)\)，由 Descartes 法则知其恰有一个正实根，记为 \(u_R\)。
最后由根的连续性与不等式在 \(u=1\) 处为真，得 \(\textsf{RH}(u)\) 在 \((0,u_R]\) 成立且在 \((u_R,\infty)\) 失效。证毕。
\end{proof}

\begin{corollary}[素轨道计数误差指数的相变：\texorpdfstring{$\beta(u)$}{beta(u)} 从 \texorpdfstring{$\le 1/2$}{<=1/2} 连续劣化到 \texorpdfstring{$1$}{1}]\label{cor:real-input-40-collision-error-phase}
在碰撞位势 \(W(u)\) 上，令 \(\lambda(u):=\lambda_1(u)=\lambda_+(u)\)，并记 \(\Lambda(u)\) 为命题 \ref{prop:real-input-40-collision-rhk-window} 的次谱半径。
则 prime-orbit 计数的主项为 \(\lambda(u)^n/n\)，误差指数由
\[
\beta(u):=\frac{\log \Lambda(u)}{\log \lambda(u)}
\]
控制：
\begin{itemize}
  \item \(0<u\le 1\) 时 \(\Lambda(u)=1\)，故 \(\beta(u)=0\)（误差至多常数量级）；
  \item \(1<u\le u_R\) 时 \(\Lambda(u)\le \sqrt{\lambda(u)}\)，故 \(\beta(u)\le \tfrac12\)（平方根相；推论 \ref{cor:kernel-rh-sqrt-error}）；
  \item \(u>u_R\) 时 \(\Lambda(u)>\sqrt{\lambda(u)}\)，故 \(\beta(u)>\tfrac12\)，并且当 \(u\to\infty\) 时 \(\beta(u)\to 1\)。
\end{itemize}
更具体地，当 \(u\to\infty\) 时三次方程的 Puiseux 展开给出
\[
\lambda_+(u)=\sqrt u+\frac12+\frac{9}{8\sqrt u}+O(u^{-1}),\qquad
\lambda_-(u)=-\sqrt u+\frac12-\frac{9}{8\sqrt u}+O(u^{-1}),
\]
从而 \(\Lambda(u)/\lambda_+(u)\to 1\)（谱隙塌缩）。
\end{corollary}
